\documentclass[11pt]{article}
% preamble.tex --- shared packages, macros, and theorem environments.
% % preamble.tex --- shared packages, macros, and theorem environments.
% % preamble.tex --- shared packages, macros, and theorem environments.
% \input{preamble} after \documentclass in each standalone note.

\usepackage{amsmath, amssymb, amsthm}
\usepackage{mathtools}
\usepackage{geometry}
\usepackage{hyperref}
\usepackage{natbib}
\usepackage{graphicx}
\usepackage{booktabs}
\usepackage{algorithm}
\usepackage{algpseudocode}
\usepackage{xcolor}

\geometry{margin=1in}

\newtheorem{proposition}{Proposition}
\newtheorem{remark}{Remark}

\DeclareMathOperator{\tr}{tr}
\DeclareMathOperator{\diag}{diag}
\newcommand{\R}{\mathbb{R}}
\newcommand{\E}{\mathbb{E}}
\newcommand{\Var}{\mathrm{Var}}
\newcommand{\dd}{\,\mathrm{d}}
\newcommand{\bx}{\mathbf{x}}
\newcommand{\bs}{\mathbf{s}}
\newcommand{\bphi}{\boldsymbol{\phi}}
 after \documentclass in each standalone note.

\usepackage{amsmath, amssymb, amsthm}
\usepackage{mathtools}
\usepackage{geometry}
\usepackage{hyperref}
\usepackage{natbib}
\usepackage{graphicx}
\usepackage{booktabs}
\usepackage{algorithm}
\usepackage{algpseudocode}
\usepackage{xcolor}

\geometry{margin=1in}

\newtheorem{proposition}{Proposition}
\newtheorem{remark}{Remark}

\DeclareMathOperator{\tr}{tr}
\DeclareMathOperator{\diag}{diag}
\newcommand{\R}{\mathbb{R}}
\newcommand{\E}{\mathbb{E}}
\newcommand{\Var}{\mathrm{Var}}
\newcommand{\dd}{\,\mathrm{d}}
\newcommand{\bx}{\mathbf{x}}
\newcommand{\bs}{\mathbf{s}}
\newcommand{\bphi}{\boldsymbol{\phi}}
 after \documentclass in each standalone note.

\usepackage{amsmath, amssymb, amsthm}
\usepackage{mathtools}
\usepackage{geometry}
\usepackage{hyperref}
\usepackage{natbib}
\usepackage{graphicx}
\usepackage{booktabs}
\usepackage{algorithm}
\usepackage{algpseudocode}
\usepackage{xcolor}

\geometry{margin=1in}

\newtheorem{proposition}{Proposition}
\newtheorem{remark}{Remark}

\DeclareMathOperator{\tr}{tr}
\DeclareMathOperator{\diag}{diag}
\newcommand{\R}{\mathbb{R}}
\newcommand{\E}{\mathbb{E}}
\newcommand{\Var}{\mathrm{Var}}
\newcommand{\dd}{\,\mathrm{d}}
\newcommand{\bx}{\mathbf{x}}
\newcommand{\bs}{\mathbf{s}}
\newcommand{\bphi}{\boldsymbol{\phi}}


\title{Sampling Methods: Reversible-Jump MCMC and\\
Infinite-Dimensional (Dimension-Robust) Samplers}
\author{Jack Heinzel}
\date{Draft --- \today}

\begin{document}
\maketitle

\begin{abstract}
This note covers the Monte Carlo machinery used to sample the discretised LGCP
posterior set up in the companion \emph{Foundations} note.  Two ingredients are
developed.  First, reversible-jump MCMC \citep{green1995} over the refinement
tree, which makes the number of active basis functions a random variable; the
birth/death moves and their acceptance ratio are presented on the triangulation
baseline and carry over verbatim to the adaptive tree bases of the companion
\emph{Universal Function Approximators} note.  Second, \emph{dimension-robust}
within-model samplers --- preconditioned Crank--Nicolson (pCN) and Hilbert-space
HMC --- whose efficiency does not degrade as the discretisation is refined,
together with a scheme that nests them inside the transdimensional moves.  The
implementation targets BlackJAX-compatible kernels.
\end{abstract}

\tableofcontents
\bigskip

% ============================================================
\section{Adaptive Refinement via Reversible-Jump MCMC}
\label{sec:rjmcmc}
% ============================================================

The number of active basis functions is itself unknown, so we sample it with
reversible-jump MCMC \citep{green1995}.  For the adaptive tree basis (companion
\emph{Universal Function Approximators} note) a \emph{birth move} advances one
node one step in its local hierarchy (const $\to$ linear $\to$ unit hat $\to$
sub-hat), adding one coefficient; a \emph{death move} reverses it.  Because
sibling hats have zero mass and stiffness overlap (the recursive-block structure
established there), each move updates the FEM matrices in $O(D)$ work with no
global rebuild.  We first present the scheme in its original form on the
triangulation baseline, where a node corresponds to a triangle and a birth
performs a ternary simplex split; the acceptance ratio below is identical for the
tree basis, with the local FEM updates replaced by the closed forms of that note.

\subsection{Refinement Tree}

We augment the static triangulation with a \emph{ternary refinement tree}.
Every triangle in the initial mesh is a root.  A \emph{birth move} splits a
leaf triangle $\tau^*$ at its barycenter $s^*$, replacing $\tau^*$ by the
three child triangles $(s^*,v_0,v_1)$, $(s^*,v_1,v_2)$, $(s^*,v_0,v_2)$.
A \emph{death move} is the exact reverse: if all three children of an
interior node $s^*$ are currently leaves, they are merged back into their
parent $\tau^*$ and $s^*$ is removed.  Birth and death are exact inverses,
so detailed balance is achievable.

Let $N_m$ be the number of leaf (birth-eligible) triangles and $M_m$ the
number of removable (death-eligible) nodes in mesh $m$.  After a birth on
$\tau^*$:
\begin{equation}
  N_{m'} = N_m + 2, \qquad
  M_{m'} = M_m + 1
           - \mathbf{1}\bigl[\text{all-leaf children test fails for parent}(s^*)\bigr],
  \label{eq:NM-update}
\end{equation}
where parent$(s^*)$ is the node whose earlier birth created $\tau^*$
(undefined for root triangles, so the indicator is 0 there).

\subsection{Conditional Prior Proposal}
\label{sec:rj-proposal}

After the 3-split the extended mesh $m'$ has $n_m+1$ nodes.  The SPDE prior
on the augmented field $(\mathbf{x}_m,x^*)$ is
$\mathcal{N}(0,\sigma^2 Q'^{-1})$ with $Q'=K'(\tilde{C}')^{-1}K'$.
Partitioning $Q'$ into the old nodes $(-*)$ and the new node $(*)$:
\begin{equation}
  Q' = \begin{pmatrix} Q'_{-*} & (Q'_{*-})^T \\ Q'_{*-} & Q'_{**} \end{pmatrix},
\end{equation}
the conditional prior of $x^*$ given the existing field is the
one-dimensional Gaussian
\begin{equation}
  x^*\mid\mathbf{x}_m,m',\kappa,\sigma
  \;\sim\; \mathcal{N}\!\left(\bar{x}^*,\;\frac{\sigma^2}{Q'_{**}}\right),
  \qquad
  \bar{x}^* = -\frac{Q'_{*-}\,\mathbf{x}_m}{Q'_{**}}.
  \label{eq:cond-prior-rj}
\end{equation}
Since $s^*$ connects only to $\{v_0,v_1,v_2\}$ in the mesh graph,
$K'_{*j}=0$ for $j\notin\{*,v_0,v_1,v_2\}$, and therefore
$Q'_{*j} = \sum_{k\in\{*,v_0,v_1,v_2\}} K'_{*k}(\tilde{c}'_k)^{-1}K'_{kj}$.
Both $Q'_{**}$ and $Q'_{*-}$ depend on $O(1)$ FEM values assembled from the
three new triangles; their evaluation is an $O(1)$ local computation.

\subsection{RJMCMC Acceptance Ratio}
\label{sec:rj-acceptance}

We propose $x^*$ from the conditional prior $q_*(x^*)=p(x^*\mid\mathbf{x}_m,m',\kappa,\sigma)$.
The standard \citet{green1995} acceptance ratio for the birth move is
\begin{equation}
  \alpha_\mathrm{b} = \min\!\left(1,\;
    \frac{\pi(m',\,\mathbf{x}_{m'})}
         {\pi(m,\,\mathbf{x}_m)}\cdot
    \frac{q_d\,N_m}{q_b\,M_{m'}\,q_*(x^*)}
  \right),
  \label{eq:rj-raw}
\end{equation}
where $q_b,q_d$ are the probabilities of attempting a birth or death move and
$\pi(m,\mathbf{x}_m)\propto p(m)\,p(\mathbf{x}_m\mid m,\kappa,\sigma)\,L(m,\mathbf{x}_m,\mu)$
is the unnormalised posterior.  (The Jacobian of the $u\to x^*$ transformation
cancels exactly against the proposal density by the change-of-variables
identity $q_*(x^*)\,|dx^*/du|^{-1}=\phi(u)$.)

\paragraph{Simplification via conditional prior proposal.}
Substituting $q_*(x^*)=p(\mathbf{x}_m,x^*\mid m')/p(\mathbf{x}_m\mid m')$
and cancelling common factors:
\begin{equation}
  \frac{p(\mathbf{x}_m,x^*\mid m')}
       {q_*(x^*)\,p(\mathbf{x}_m\mid m)}
  = \frac{p(\mathbf{x}_m\mid m')}{p(\mathbf{x}_m\mid m)}
  \;\eqqcolon\; \mathcal{R}_\mathrm{prior},
\end{equation}
so the acceptance ratio reduces to
\begin{equation}
  \alpha_\mathrm{b} = \min\!\left(1,\;
    \mathcal{R}_\mathrm{prior}
    \cdot
    \underbrace{\frac{L(m',\mathbf{x}_{m'},\mu)}
                     {L(m,\mathbf{x}_m,\mu)}}_{\mathcal{R}_L}
    \cdot
    \underbrace{\frac{p(m')}{p(m)}}_{\mathcal{R}_m}
    \cdot\frac{q_d\,N_m}{q_b\,M_{m'}}
  \right).
  \label{eq:rj-acc}
\end{equation}
The death move acceptance ratio is $\min(1,\alpha_\mathrm{b}^{-1})$.

\subsection{Prior Consistency and the Approximation $\mathcal{R}_\mathrm{prior}\approx 1$}
\label{sec:rj-prior}

The term $\mathcal{R}_\mathrm{prior} = p(\mathbf{x}_m\mid m')/p(\mathbf{x}_m\mid m)$
compares the SPDE prior on the \emph{old} nodes under the two meshes.
For a true Mat\'ern Gaussian process, this ratio equals 1 exactly: the GP
consistency property guarantees that the marginal at any fixed set of
locations is determined solely by the covariance function, independent of
which additional locations are included in the model.

For the \emph{FEM-discretised} prior, however, $\mathcal{R}_\mathrm{prior}\neq 1$
in general.  Splitting $\tau^*$ modifies the piecewise-linear basis functions
at $v_0,v_1,v_2$: the stiffness and mass integrals over the former $\tau^*$
are replaced by integrals over three smaller triangles, so $Q'_{-*}\neq Q$ and
the Schur complement $Q'_\mathrm{SC} \neq Q$.  The discrepancy is
\begin{equation}
  Q'_\mathrm{SC} - Q \;=\;
  \underbrace{(Q'_{-*} - Q)}_{\text{FEM change at }v_0,v_1,v_2}
  - \frac{(Q'_{*-})^T Q'_{*-}}{Q'_{**}},
  \label{eq:SC-diff}
\end{equation}
which has nonzeros only in rows/columns $\{v_0,v_1,v_2\}$ and their
2-hop neighbours.  Both terms scale with the area $|\tau^*|$ of the split
triangle, so $\|Q'_\mathrm{SC}-Q\|=O(|\tau^*|)$ and
$\log\mathcal{R}_\mathrm{prior}=O(|\tau^*|^2)$.  Finer triangles produce
smaller corrections, and the correction vanishes in the mesh-refinement limit
as the FEM prior converges to the true Matérn GP.

We therefore adopt the approximation
\begin{equation}
  \mathcal{R}_\mathrm{prior} \approx 1,
  \label{eq:Rprior-approx}
\end{equation}
which is exact for the true GP prior and is a consistent approximation for
the FEM prior.  The acceptance ratio simplifies to
\begin{equation}
  \boxed{
  \alpha_\mathrm{b} = \min\!\left(1,\;
    \underbrace{\frac{L(m',\mathbf{x}_{m'},\mu)}
                     {L(m,\mathbf{x}_m,\mu)}}_{\mathcal{R}_L}
    \cdot
    \underbrace{\frac{p(m')}{p(m)}}_{\mathcal{R}_m}
    \cdot\frac{q_d\,N_m}{q_b\,M_{m'}}
  \right).}
  \label{eq:rj-acc-simple}
\end{equation}

\begin{remark}[Exact cancellation via GP prior]
$\mathcal{R}_\mathrm{prior}=1$ holds \emph{exactly} if the prior on every mesh
is defined by the Mat\'ern covariance function directly---i.e.\
$p(\mathbf{x}_m\mid m,\kappa,\sigma)=\mathcal{N}(0,\Sigma_m)$ with
$\Sigma_{ij}=\sigma^2 C_\nu(s_i,s_j;\kappa)$---rather than by the FEM
precision.  In that case the proposal \eqref{eq:cond-prior-rj} becomes the
kriging predictor under the Mat\'ern covariance, replacing the local FEM
formula with a dense covariance solve.  The FEM approximation is then used
only for the likelihood (lumped-mass quadrature and barycentric interpolation),
not for the prior.
\end{remark}

\subsection{Likelihood Ratio}
\label{sec:rj-ll}

Because only the triangles touching $\tau^*$ change, the log-likelihood ratio
is local:
\begin{equation}
  \log\mathcal{R}_L
  = \!\sum_{i:\,s_i\in\tau^*}\!\!\!\log\frac{\hat\lambda'(s_i)}{\hat\lambda(s_i)}
  - \left[\tilde{c}'_{s^*}e^{\mu+x^*}
    + \sum_{j\in\{v_0,v_1,v_2\}}\!(\tilde{c}'_{v_j}-\tilde{c}_{v_j})\,e^{\mu+x_{v_j}}\right].
  \label{eq:rj-ll}
\end{equation}
The first sum covers only the $O(n/N_m)$ observations falling inside $\tau^*$,
using barycentric weights recomputed in the new child triangles.

\subsection{Mesh Prior}
\label{sec:mesh-prior}

A natural choice is a geometric prior on the number of leaf triangles:
\begin{equation}
  p(m) \propto \rho^{N_m-N_0}, \qquad 0<\rho<1,
  \label{eq:mesh-prior}
\end{equation}
where $N_0$ is the initial triangle count and $\rho$ controls the expected
depth of refinement.  Under \eqref{eq:mesh-prior}, $\mathcal{R}_m=\rho^2$
for every birth move (two new leaves added).
% ============================================================
\section{Infinite-Dimensional Sampling Methods}
\label{sec:inf-dim-samplers}
% ============================================================

The RJMCMC scheme of Section~\ref{sec:rjmcmc} handles the \emph{transdimensional}
part of the posterior — deciding how many mesh nodes to use.  Within a fixed mesh,
however, we still need to sample the field values $x \in \mathbb{R}^N$.  This
section explains why the obvious choices (random-walk MH, vanilla HMC) become
increasingly expensive as $N$ grows, and introduces two \emph{dimension-robust}
alternatives that are designed from the start to work on function space.

\subsection{The Curse of Discretisation}
\label{sec:curse-disc}

Fix a prior $x \sim \mathcal{N}(0, Q^{-1})$ with precision $Q \in \mathbb{R}^{N\times N}$.
As the mesh is refined, $N\to\infty$, and the field $x$ converges to a genuine
Gaussian process.  Standard samplers degrade:

\begin{itemize}
  \item \textbf{Random-walk MH} with isotropic proposal variance $h^2 I$:
        the optimal acceptance rate requires $h = O(N^{-1/2})$, so the chain
        makes $O(N)$ steps to move $O(1)$ in function space — cost $O(N)$ per
        independent sample.
  \item \textbf{Vanilla HMC} (mass matrix $M = I$): the leapfrog step size must
        satisfy $\varepsilon = O(N^{-1/4})$ to keep $|\Delta H|$ bounded, so
        $O(N^{1/4})$ steps per proposal and $O(N^{1/4})$ gradient evaluations,
        giving $O(N^{5/4})$ cost per independent sample.
\end{itemize}

The root cause is the same in both cases: the proposal ignores the prior geometry,
so the acceptance probability collapses in high dimensions.  The fix is to build
the prior covariance $Q^{-1}$ directly into the proposal.

\subsection{Preconditioned Crank--Nicolson (pCN)}
\label{sec:pcn}

\paragraph{Proposal.}
For a target $\pi(x) \propto L(x)\,\mathcal{N}(x;\,0,Q^{-1})$, define:
\begin{equation}
  x' = \sqrt{1-\beta^2}\,x + \beta\,\xi, \qquad \xi \sim \mathcal{N}(0, Q^{-1}),
  \label{eq:pcn-proposal}
\end{equation}
where $\beta \in (0,1]$ is a step-size parameter.

\paragraph{Why it works.}
The Crank--Nicolson interpolation \eqref{eq:pcn-proposal} exactly preserves
$\mathcal{N}(0, Q^{-1})$: if $x$ is distributed as the prior, so is $x'$.
Consequently, the prior log-density cancels in the Metropolis--Hastings ratio and
only the likelihood ratio survives:
\begin{equation}
  \alpha(x, x') = \min\!\left(1,\, \frac{L(x')}{L(x)}\right).
  \label{eq:pcn-alpha}
\end{equation}
The acceptance rate therefore depends only on $L$, not on $N$: choosing $\beta$
once for a coarse mesh yields the same acceptance on a fine mesh.  This
\emph{dimension-robustness} is the defining property of pCN
\citep{cotter2013}.

\paragraph{Implementation.}
Sampling $\xi \sim \mathcal{N}(0, Q^{-1})$ amounts to solving the sparse linear
system $Q\,\xi = z$ for $z \sim \mathcal{N}(0,I)$, which can be done with
conjugate gradient (CG) in $O(N)$ time for a graph-Laplacian precision (see
Section~\ref{sec:cg}).

\subsection{Hilbert-Space HMC}
\label{sec:hilbert-hmc}

pCN is robust but uses only first-order (likelihood evaluation) information.
Hilbert-space HMC incorporates gradient information while retaining
dimension-robustness \citep{beskos2011}.

\paragraph{Velocity parameterisation.}
Standard HMC introduces a momentum $p \sim \mathcal{N}(0, M)$ and
Hamiltonians $H = U(x) + \tfrac{1}{2}p^\top M^{-1} p$.  Setting $M = Q$
and changing variables to $v = Q^{-1}p$ (so $v \sim \mathcal{N}(0, Q^{-1})$):
\begin{align}
  H(x,v) &= U(x) + K(v), \\
  U(x)   &= -\log L(x) + \tfrac{1}{2} x^\top Q x, \label{eq:hmc-U}\\
  K(v)   &= \tfrac{1}{2} v^\top Q v. \label{eq:hmc-K}
\end{align}
Hamilton's equations become $\dot{x} = v$ and $Q\dot{v} = -\nabla U(x)$.

\paragraph{Leapfrog integrator.}
With step size $\varepsilon$ and $L$ leapfrog steps:
\begin{align}
  v &\leftarrow v - \tfrac{\varepsilon}{2}\, Q^{-1}\nabla U(x) \tag{half kick}\\
  \text{for } & i = 1,\ldots, L-1: \notag\\
  & x \leftarrow x + \varepsilon\, v \tag{drift}\\
  & v \leftarrow v - \varepsilon\, Q^{-1}\nabla U(x) \tag{full kick}\\
  x &\leftarrow x + \varepsilon\, v \tag{final drift}\\
  v &\leftarrow v - \tfrac{\varepsilon}{2}\, Q^{-1}\nabla U(x). \tag{final half kick}
\end{align}
Each ``kick'' requires one solve $Q^{-1}b$; each drift requires only a
vector addition.  The kinetic energy $K(v) = \tfrac{1}{2}v^\top Qv$ needs only
a matrix–vector product $Qv$, not a solve.

\paragraph{Why the velocity form is preferable.}
Three advantages over the standard $(x, p)$ form with $M = Q$:
\begin{enumerate}
  \item $K(v) = \tfrac{1}{2}v^\top Qv$ — computed with one sparse mat-vec, no
        matrix inverse.
  \item Velocity refresh $v \sim \mathcal{N}(0, Q^{-1})$ uses \emph{exactly} the same
        CG solve as prior sampling — no Cholesky, no dense matrix stored.
  \item All $O(N^2)$--$O(N^3)$ dense-algebra is replaced by $O(N \times
        \text{CG iters})$ sparse work.
\end{enumerate}

\paragraph{Dimension-robustness.}
With $M = Q$, the leapfrog step size can be chosen independently of $N$ (the
eigenmodes of $Q$ all enter the dynamics on the same footing).  In practice
one observes $|\Delta H| = O(1)$ and acceptance $\approx 65$--$80\,\%$ regardless
of mesh size, provided $\varepsilon$ is tuned on a coarse mesh \citep{beskos2011,
hairer2014}.

\paragraph{CG for $Q^{-1}b$.}
\label{sec:cg}
For the Besag CAR prior $Q = \kappa^2 I + L_{\rm graph}$, where $L_{\rm graph}$
is the graph Laplacian of the mesh connectivity, the condition number satisfies
\begin{equation}
  \kappa(Q) = \frac{\kappa^2 + \lambda_{\max}(L_{\rm graph})}{\kappa^2}
            \approx \frac{\kappa^2 + 8}{\kappa^2}
\end{equation}
for a 4-connected grid (eigenvalues of $L_{\rm graph}$ lie in $[0,8]$).  This is
\emph{independent of $N$}, so CG converges in $O(1)$ iterations and each leapfrog
step costs $O(N)$.  For the FEM precision matrix derived from the SPDE operator
(companion \emph{Foundations} note), the same argument holds with $\kappa^2$
replacing 8 as the dominant eigenvalue bound.

\subsection{Combining Hilbert-Space HMC with RJMCMC}
\label{sec:rjmcmc-hmc}

The overall algorithm separates into two nested move types:

\begin{description}
  \item[\textbf{Within-model move}] (fixed $N_m$): one Hilbert-space HMC step in
    $(x, v)$ space.  Gradient information from the likelihood accelerates mixing;
    the $Q$-preconditioned leapfrog keeps $\varepsilon$ mesh-size-independent.
  \item[\textbf{Transdimensional move}] (birth or death): one RJMCMC step as in
    Section~\ref{sec:rjmcmc}.  On acceptance, the field $x$ and precision $Q$
    change dimension.  The velocity $v$ is refreshed by drawing
    $v_{\rm new} \sim \mathcal{N}(0, Q_{\rm new}^{-1})$
    (one CG solve in the new space) before the next HMC step.
\end{description}

\begin{algorithm}[H]
\caption{Joint RJMCMC + Hilbert-Space HMC}
\label{alg:rj-hmc}
\begin{algorithmic}[1]
  \State \textbf{Initialise:} mesh $\mathcal{T}_0$, $x^{(0)} \sim \mathcal{N}(0, Q_0^{-1})$,
         $v^{(0)} \sim \mathcal{N}(0, Q_0^{-1})$
  \For{$t = 1, 2, \ldots$}
    \State Draw $u \sim \mathrm{Uniform}(0,1)$
    \If{$u < p_{\rm trans}$}  \Comment{transdimensional move}
      \State Propose birth/death $(\mathcal{T}, x) \to (\mathcal{T}', x')$ \hfill (Section~\ref{sec:rjmcmc})
      \State Accept with probability $\min(1,\, \mathcal{R}_L \cdot \mathcal{R}_m \cdot \mathcal{R}_J)$
      \If{accepted}
        \State $v \leftarrow Q'^{-1} z$, \;$z \sim \mathcal{N}(0, I)$
        \hfill (velocity refresh via CG)
      \EndIf
    \Else \Comment{within-model move}
      \State $v \leftarrow Q^{-1}z$, \;$z \sim \mathcal{N}(0, I)$
            \hfill (velocity refresh)
      \State Run $L$-step leapfrog to propose $(x', v')$
      \State Accept with probability $\min(1,\, e^{-\Delta H})$
    \EndIf
  \EndFor
\end{algorithmic}
\end{algorithm}

\noindent
The mixing-time scales of the two move types are naturally separated: HMC
explores the field posterior quickly (correlated draws in $O(1)$ steps) while
RJMCMC explores mesh complexity slowly (each birth/death requires a Jacobian
factor).  The fraction $p_{\rm trans}$ controls the trade-off; a common heuristic
is $p_{\rm trans} \approx 0.1$.

\subsection{Annotated Reading List}
\label{sec:reading-list}

\begin{description}
  \item[\citet{stuart2010}] The foundational reference for Bayesian inversion on
    function space.  Section 4 defines Gaussian priors on Hilbert spaces and makes
    precise what it means for a posterior to be well-posed in infinite dimensions.
    Read this first to understand \emph{why} dimension-robustness is the right
    criterion.

  \item[\citet{cotter2013}] Derives pCN and its acceptance formula \eqref{eq:pcn-alpha},
    proves dimension-robustness, and compares it to a function-space version of
    the Langevin algorithm (MALA).  Includes numerical experiments on an elliptic
    inverse problem showing acceptance rates stable as $N \to \infty$.

  \item[\citet{beskos2011}] Introduces the Hilbert-space HMC (called ``Hybrid Monte
    Carlo on Hilbert spaces'').  Proves that with the prior as mass matrix the
    step size can be held fixed as $N \to \infty$ and that the acceptance
    probability remains bounded away from zero.  The velocity parameterisation
    used here follows their exposition.

  \item[\citet{hairer2014}] Provides spectral-gap bounds for pCN and related
    infinite-dimensional Metropolis algorithms, giving a rigorous convergence rate
    in terms of the likelihood Lipschitz constant.  Useful background if you want
    to understand \emph{how fast} the chain mixes in theory.

  \item[\citet{neal2011}] The standard handbook chapter on HMC; covers the
    leapfrog integrator, step-size tuning, and the dual-averaging adaptation used
    in NUTS.  Read before \citet{beskos2011} to understand the finite-dimensional
    version.

  \item[\citet{green1995}] Original RJMCMC paper.  Introduces the general
    reversible-jump framework and the Jacobian term $\mathcal{R}_J$; the birth/death
    moves of Section~\ref{sec:rjmcmc} are a special case.
\end{description}

\bibliographystyle{plainnat}
\bibliography{refs}

\end{document}
